\section{Group Theory}

This section focuses on the group-theoretic tools needed for the course: in particular, the finite abelian groups of units modulo an integer and their multiplicative characters. The selection follows the introductory group material in Lee \cite[Chapters 3--5]{Lee2018-aa}.

\subsection{Groups and Groups of Units}

\begin{definition}[Group]
A group is a set $G$ equipped with a binary operation, written multiplicatively, such that:
\begin{enumerate}
    \item the operation is associative;
    \item there is an identity element $e\in G$ satisfying $eg=ge=g$ for every $g\in G$;
    \item every $g\in G$ has an inverse $g^{-1}\in G$ satisfying $gg^{-1}=g^{-1}g=e$.
\end{enumerate}
A group is abelian if $gh=hg$ for all $g,h\in G$.
\end{definition}

\begin{definition}[Subgroup]
A subset $H\subseteq G$ is a subgroup, written $H\leq G$, if it is a group under the operation inherited from $G$.
\end{definition}

The group relevant to Dirichlet characters is
\[
 U(m)=(\mathbb{Z}/m\mathbb{Z})^{\times},
\]
the group of invertible residue classes modulo $m$. As multiplication of residue classes is commutative, $U(m)$ is abelian. Its order is $\varphi(m)$, the number of residue classes relatively prime to $m$.

\begin{example}
For $m=5$, $U(5)=\{[1]_5,[2]_5,[3]_5,[4]_5\}$. Since
\[
 [2]_5^2=[4]_5,\qquad [2]_5^3=[3]_5,\qquad [2]_5^4=[1]_5,
\]
the group is cyclic, generated by $[2]_5$.
\end{example}

\subsection{Cosets, Orders, and Lagrange's Theorem}

\begin{definition}[Coset]
If $H\leq G$ and $g\in G$, the left coset of $H$ determined by $g$ is
\[
 gH=\{gh:h\in H\}.
\]
\end{definition}

The left cosets of $H$ partition $G$, and each has the same number of elements as $H$.

\begin{definition}[Order of an element]
For $g\in G$, the cyclic subgroup generated by $g$ is $\langle g\rangle=\{g^k:k\in\mathbb{Z}\}$. If this subgroup is finite, its number of elements is called the order of $g$ and is denoted $|g|$.
\end{definition}

\begin{theorem}[Lagrange's theorem]
If $G$ is a finite group and $H\leq G$, then
\[
 |G|=[G:H]|H|,
\]
where $[G:H]$ is the number of distinct left cosets of $H$ in $G$. In particular, the order of every element of $G$ divides $|G|$.
\end{theorem}

\begin{proof}
The left cosets partition $G$ and the map $h\mapsto gh$ is a bijection from $H$ to $gH$. Thus every coset has $|H|$ elements, and counting the elements of $G$ by cosets gives $|G|=[G:H]|H|$. For any $g\in G$, the powers of $g$ form the cyclic subgroup $\langle g\rangle$; applying the formula to this subgroup shows that the order of $g$ divides $|G|$.
\end{proof}

\subsection{Normal Subgroups and Quotient Groups}

\begin{definition}[Normal subgroup]
A subgroup $N\leq G$ is normal, written $N\trianglelefteq G$, if $gN=Ng$ for every $g\in G$.
\end{definition}

\begin{definition}[Quotient group]
If $N\trianglelefteq G$, the set of cosets
\[
G/N=\{gN:g\in G\}
\]
is a group under the operation $(gN)(hN)=(gh)N$. Normality ensures that this operation does not depend on the choice of coset representatives. The group is called the quotient group of $G$ by $N$.
\end{definition}

\begin{example}
The subgroup $m\mathbb{Z}$ is normal in the additive group $\mathbb{Z}$, and the quotient group $\mathbb{Z}/m\mathbb{Z}$ consists of the residue classes modulo $m$.
\end{example}

\subsection{Homomorphisms and the First Isomorphism Theorem}

\begin{definition}[Group homomorphism]
Let $G$ and $H$ be groups. A function $\varphi:G\to H$ is a group homomorphism if
\[
 \varphi(g_1g_2)=\varphi(g_1)\varphi(g_2)
 \qquad\text{for all }g_1,g_2\in G.
\]
Its kernel and image are
\[
 \ker\varphi=\{g\in G:\varphi(g)=e_H\},\qquad
 \operatorname{im}\varphi=\{\varphi(g):g\in G\}.
\]
An isomorphism is a bijective homomorphism; groups related by an isomorphism have the same group structure.
\end{definition}

If $N\trianglelefteq G$ and a homomorphism $\varphi:G\to H$ satisfies $N\subseteq\ker\varphi$, then $\varphi$ is constant on each coset of $N$. It therefore induces a homomorphism $\overline{\varphi}:G/N\to H$ defined by $\overline{\varphi}(gN)=\varphi(g)$, so that $\varphi=\overline{\varphi}\circ\pi$, where $\pi(g)=gN$ is the quotient map. The diagram illustrates this factorization:

\begin{figure}[H]
\centering
\begin{tikzcd}[column sep=large,row sep=large]
G \arrow[r,"\varphi"] \arrow[d,"\pi"'] & H \\
G/N \arrow[ur,"\overline{\varphi}"']
\end{tikzcd}
\caption{A homomorphism that is trivial on $N$ factors through the quotient $G/N$.}
\label{fig:quotient-factorization}
\end{figure}

\begin{theorem}[Kernel and image]
For a group homomorphism $\varphi:G\to H$, the kernel is a normal subgroup of $G$ and the image is a subgroup of $H$.
\end{theorem}

\begin{proof}
If $h_1=\varphi(g_1)$ and $h_2=\varphi(g_2)$ are in the image, then $h_1h_2^{-1}=\varphi(g_1g_2^{-1})$ is also in the image, so the image is a subgroup. Since $\varphi(e_G)=e_H$ and $\varphi(k_1k_2^{-1})=e_H$ whenever $k_1,k_2\in\ker\varphi$, the kernel is a subgroup. If $k\in\ker\varphi$ and $g\in G$, then
\[
 \varphi(gkg^{-1})=\varphi(g)\varphi(k)\varphi(g)^{-1}
 =\varphi(g)e_H\varphi(g)^{-1}=e_H.
\]
Thus $g(\ker\varphi)g^{-1}\subseteq\ker\varphi$; applying the same argument to $g^{-1}$ gives equality, so the kernel is normal.
\end{proof}

\begin{theorem}[First isomorphism theorem]
For every group homomorphism $\varphi:G\to H$,
\[
 G/\ker\varphi\cong\operatorname{im}\varphi.
\]
\end{theorem}

The theorem identifies the quotient by the kernel with the image. The diagram shows the quotient map, the induced isomorphism, and the inclusion of the image into $H$:

\begin{figure}[H]
\centering
\begin{tikzcd}[column sep=large,row sep=large]
G \arrow[r,"\varphi"] \arrow[d,"\pi"'] & H \\
G/\ker\varphi \arrow[r,"\overline{\varphi}"'] & \operatorname{im}\varphi \arrow[u,hook,"i"']
\end{tikzcd}
\caption{The first isomorphism theorem: $\overline{\varphi}$ is an isomorphism and $i$ is the inclusion.}
\label{fig:first-isomorphism}
\end{figure}

\begin{proof}
Define $\overline{\varphi}:G/\ker\varphi\to\operatorname{im}\varphi$ by $\overline{\varphi}(g\ker\varphi)=\varphi(g)$. If two cosets are equal, their representatives differ by an element of the kernel, so they have the same image; hence the map is well-defined. It is a homomorphism, is onto by the definition of the image, and is one-to-one because equal images imply that the representatives differ by an element of the kernel. Therefore it is an isomorphism.
\end{proof}

\subsection{Finite-Group Characters}

\begin{definition}[Multiplicative character]
A multiplicative character of a group $G$ is a homomorphism $\chi:G\to\mathbb{C}^{\times}$, where $\mathbb{C}^{\times}=\mathbb{C}\setminus\{0\}$ under multiplication. For a finite abelian group, its character group is
\[
 \widehat{G}=\operatorname{Hom}(G,\mathbb{C}^{\times}).
\]
The operation on $\widehat{G}$ is pointwise multiplication of characters.
\end{definition}

\begin{proposition}
If $G$ is finite and $\chi:G\to\mathbb{C}^{\times}$ is a character, then every value $\chi(g)$ is a root of unity.
\end{proposition}

\begin{proof}
For $g\in G$, let $r$ be its finite order. Then $g^r=e$, and the homomorphism property gives $\chi(g)^r=\chi(g^r)=\chi(e)=1$.
\end{proof}

For the course, take $G=U(m)$. A multiplicative character on $U(m)$ gives the values of a Dirichlet character on the invertible residue classes modulo $m$; it is extended to integers by assigning value $0$ to integers not relatively prime to $m$. For example, since $U(5)=\langle[2]_5\rangle$, one character is determined by $\chi([2]_5)=i$, yielding $\chi([1]_5)=1$, $\chi([2]_5)=i$, $\chi([4]_5)=-1$, and $\chi([3]_5)=-i$. These are the group-theoretic objects used to define Dirichlet $L$-functions.